% Generated draft: requires mathematical and source review.
\documentclass[11pt]{article}
\usepackage[margin=25mm]{geometry}
\usepackage{fontspec}
\setmainfont{Latin Modern Roman}
\usepackage{amsmath,amssymb,mathtools,unicode-math}
\setmathfont{Latin Modern Math}
\usepackage{enumitem,needspace,xurl,hyperref}
\setlength{\parindent}{0pt}
\setlength{\parskip}{5pt}
\setlength{\emergencystretch}{3em}
\newcommand{\N}{\mathbb{N}}
\newcommand{\Z}{\mathbb{Z}}
\newcommand{\Q}{\mathbb{Q}}
\newcommand{\R}{\mathbb{R}}
\newcommand{\C}{\mathbb{C}}
\newcommand{\F}{\mathbb{F}}
\newcommand{\A}{\mathbb{A}}
\newcommand{\PP}{\mathbb{P}}
\newcommand{\E}{\mathbb{E}}
\newcommand{\eps}{\varepsilon}
\newcommand{\dd}{\,\mathrm{d}}
\begin{document}
\section*{UnsolvedMath ID: 30006097}
\textbf{Source problem code:} catalogue entry
\subsection{Definitions and mathematical statement}
A reduced-order model is a model whose state space has smaller dimension than that of the original system. When these dimensions are finite, denote them by \(r\) and \(n\), respectively, so that \(r<n\). The supplied record concerns balanced reduction for nonlinear input--output systems, but does not define the balancing construction, the associated energy functions, or the Hamilton--Jacobi--Bellman equations to be approximated. Output-feedback control means that control uses observed outputs; the available observations and permitted controller dynamics are not specified.

\textbf{Research problem as supplied.} Develop efficient balanced reduced-order models for nonlinear systems and compare them systematically with other nonlinear model-reduction methods. In particular:
\begin{enumerate}
\item Develop scalable approximations of the associated Hamilton--Jacobi--Bellman equations that enable real-time output-feedback control.
\item Determine whether iterative solvers and alternative basis functions can replace direct polynomial-coefficient solvers, reducing both runtime and memory sufficiently to handle substantially higher state-space dimensions.
\item Investigate the nonlinear transformation approach mentioned in the original excerpt, described there as \emph{balance-and-reduce}, in contrast to \emph{balance-then-reduce}.
\end{enumerate}

\textbf{Unresolved formulation qualifications.} These requests constitute a research programme rather than a specified conjecture or estimate. The record supplies no class of nonlinear systems, state or input domains, regularity assumptions, initial or boundary conditions, Hamilton--Jacobi--Bellman equations, solution concept, approximation tolerance, comparison protocol, or computational resource model. It gives no quantitative meanings for ``efficient,'' ``scalable,'' ``real-time,'' or ``substantially higher.'' The original excerpt is fragmentary and does not specify the transformation approach or its relationship to the expanded statement. Consequently, no definite universally or existentially quantified mathematical assertion can be extracted without additional source information.
\subsection{Short English statement}
Develop and compare computationally efficient nonlinear balanced model-reduction methods, including scalable Hamilton--Jacobi--Bellman approximations for output-feedback control and alternatives to direct polynomial-coefficient solvers. The supplied material leaves the system class, equations, comparison criteria, and quantitative targets unspecified, so this is a statement-only research-programme draft requiring mathematical review.
\subsection{Sources}
\begin{itemize}
\item Local UnsolvedMath registry, supplied complete record for permanent ID \(30006097\), title \emph{Efficient Nonlinear Balanced-Truncation Model Reduction}, exact source problem code \texttt{catalogue entry}, in the collection \emph{Oberwolfach Reports Open Problems}. No individual proposer is supplied.
\item Supplied source citation: \emph{Mini-Workshop: Data-driven Modeling, Analysis, and Control of Dynamical Systems (2025)}, Oberwolfach Reports, pp.\ 3255--3300, DOI \texttt{10.4171/owr/2024/57}. Supplied source URL: \url{https://doi.org/10.4171/owr/2024/57}. The background describes Oberwolfach Reports as records of mathematical workshops and their open-problem discussions.
\item Supplied literature source: \emph{Oberwolfach Report 57}, \url{https://doi.org/10.4171/OWR/2024/57}.
\item Literature assessment dated exactly \texttt{2026-08-22}: the supplied curation labels the problem open, says that the workshop explicitly leaves open problems for nonlinear balanced truncation, and observes that the extracted record does not specify a single resolved conjecture or estimate. This is supplied dated curation, not an independently verified literature assessment; its time-sensitive status and cited claims require independent verification.
\item The registry supplies proposed year \(2024\) and tag \texttt{owr-year:2024}, while the source citation names the workshop with year \(2025\). This discrepancy is unresolved. The supplied statement status is \texttt{corrected\_verified}; its accompanying verification note says that only clear formula delimiters, notation, grammar, and typography were repaired and that mathematical content and scope were preserved. These are source-record assertions, not independent verification of this draft.
\end{itemize}
\end{document}
